\documentclass[12pt,a4paper]{article}

\usepackage[a4paper,left=2.0cm,right=2.0cm,top=2.8cm,bottom=2.5cm]{geometry}
\usepackage{fontspec}
\usepackage{enumitem}
\usepackage{microtype}
\usepackage{setspace}
\usepackage{titlesec}
\usepackage{amsmath,amssymb}

\setmainfont{Times New Roman}
\setstretch{1.3}
\setlength{\parindent}{0pt}
\setlength{\parskip}{0.55em}
\setlist[itemize]{leftmargin=2em,itemsep=0.35em,topsep=0.35em}
\setlist[itemize,2]{label=--,leftmargin=2.2em}

\titleformat{\section}{\Large\bfseries}{\thesection}{1em}{}
\titleformat{\subsection}{\large\bfseries}{\thesubsection}{1em}{}
\titlespacing*{\section}{0pt}{1.1em}{0.45em}
\titlespacing*{\subsection}{0pt}{0.9em}{0.35em}

\begin{document}

\begin{center}
  {\LARGE Machine Learning Course 2026 Autumn: Experiment 2}

  \vspace{1.35cm}
  {\large October 13, 2026}
\end{center}

\vspace{0.7cm}

\section{Overview}

In this experiment, you will implement linear and nonlinear Support Vector
Machine (SVM) models on given datasets, and deepen your understanding of the
hyperparameters and kernel functions through visualization. You will also
implement a linear SVM using \textbf{NumPy} and compare it with the model in
\textbf{\textit{sklearn}}. Tasks 1--4 are required; Task 5 is optional.

\section{Dataset Description}

\begin{sloppypar}
The datasets are ready for you and are stored in the files
\textbf{\textit{T1\_data.csv}} and \textbf{\textit{T2\_data.csv}}.
Each row is a sample, with the last column being the label and the other columns
being the features. T1 has 51 samples and T2 has 200 samples; both have two
features and labels $-1$ and $+1$. You also need to load the
\textbf{Wisconsin Breast Cancer Dataset} using \textbf{\textit{sklearn}}:
569 samples, 30 features, and original labels 0 (malignant) and 1 (benign).
\end{sloppypar}

\section{Task Description}

\subsection{Task 1}

Given the \textbf{\textit{T1}} dataset, fit a \textbf{linear SVM} with $C=1$.
Visualize the data, decision boundary, margin contours and support vectors.
Compare $C\in\{0.01,0.1,1,10,100,1000\}$, and explain how $C$ changes the
margin and the penalty for margin violations. These full-dataset plots illustrate
the fitted models; training accuracy does not estimate unseen-data performance.

\subsection{Task 2}

\begin{itemize}
  \item Given the \textbf{\textit{T2}} dataset, fit a \textbf{nonlinear SVM with
        Gaussian (RBF) kernel}, initially with $C=1$ and $\gamma=1$.
        Visualize the nine combinations of $C\in\{0.1,1,10\}$ and
        $\gamma\in\{0.1,1,10\}$. Explain their different effects on the
        decision boundary. The kernel is
        \[
          K(x,z)=\exp\!\left(-\gamma\lVert x-z\rVert^2\right),
          \qquad \gamma=\frac{1}{2\sigma^2}.
        \]
        Larger $\gamma$ corresponds to a narrower Gaussian kernel. In the
        soft-margin dual, the constraints are $0\leq\alpha_i\leq C$ and
        $\sum_i\alpha_i y_i=0$.

  \item Tune $C$ and $\gamma$ using a \textbf{grid search}. Use all 21 values
        $2^{-10},2^{-9},\ldots,2^{10}$ for each parameter, giving 441 candidates.
        Use \textbf{5-fold stratified cross-validation}, with shuffle enabled,
        random state 42 and accuracy scoring. Report the optimal parameters and
        mean validation accuracy. Plot the 21-by-21 score heatmap with labelled
        logarithmic parameter axes. The best CV score is a tuning result, not an
        independent test score. Other kernel comparisons are optional.
\end{itemize}

\subsection{Task 3}

\begin{itemize}
  \item Load the \textbf{Wisconsin Breast Cancer Dataset} as a feature DataFrame
        and target Series. Split it into a \textbf{training set (80\%)} and a
        \textbf{test set (20\%)} with stratification and random state 42.
        Keep raw features in \texttt{X\_train} and \texttt{X\_test}.
        Demonstrate standardization on training data in a separate variable and
        explain why feature scaling matters for SVMs.

  \item For model selection, use a \textbf{Pipeline} consisting of
        \textbf{StandardScaler} and an \textbf{RBF SVC}. The scaler must be fitted
        inside each training fold, rather than before cross-validation.
        Search exactly seven values $10^{-3},10^{-2},\ldots,10^3$ for $C$ and
        seven values $10^{-4},10^{-3},\ldots,10^2$ for $\gamma$.
        Use shuffled, stratified five-fold CV, random state 42, and
        \textbf{malignant F1} scoring with \texttt{pos\_label=0}.
        Fit the search only on the raw training subset.

  \item After selecting the parameters using CV, evaluate the best Pipeline once
        on the test set. Report \textbf{accuracy, precision, recall and F1}.
        Retain the original labels: \textbf{0 = malignant, 1 = benign}; malignant
        is the positive class for the binary metrics. Include both class reports
        and a labelled confusion matrix in order $[0,1]$. Explain the malignant
        false negatives (true malignant predicted benign) and the difference
        between tuning CV and independent test performance.
\end{itemize}

\subsection{Task 4}

Use the provided data preparation, model interface and training-loop framework
to implement a \textbf{linear SVM using NumPy}. The objective is
\[
  L(w,b)=\frac12\lVert w\rVert^2+
  C\sum_i\max\!\left(0,1-y_i(w^T x_i+b)\right).
\]
The hinge penalties are \textbf{summed}, not averaged, and the bias is not
regularized.

\begin{itemize}
  \item Reload T1 and standardize it with NumPy means and population standard
        deviations, replacing any zero standard deviation by 1. Complete the
        loss, subgradients, parameter updates and prediction functions in
        \texttt{NumPyLinearSVM}. Provide \texttt{fit},
        \texttt{decision\_function} and \texttt{predict}; predict $+1$ when the
        decision score is zero. Derive both subgradients in your report.

  \item Initialize weights and bias to zero. Perform \textbf{5000 full-batch
        updates}. Use subgradients and the learning rate
        $\eta_t=0.05/\sqrt{t+1}$ for
        $t=0,\ldots,4999$. Samples with $y_i(w^T x_i+b)<1$ are active;
        choose the inactive branch at exactly margin 1. Save the training loss
        history and the parameters with the lowest objective, including the
        initial state. Run the supplied numerical gradient and score-zero checks.

  \item For $C\in\{0.1,1,10\}$, compare with \texttt{SVC} using
        \texttt{kernel="linear"}, the same $C$, and \texttt{tol=1e-8}.
        Use the same standardized data for both models. Report training accuracy,
        objective value and margin width
        $2/\lVert w\rVert$ in standardized coordinates. Losses must be finite,
        the best loss must improve on initialization, and the relative objective
        difference from sklearn must be at most 5\%.

  \item Plot the raw and best-so-far loss curves and six decision boundaries.
        Analyze the differences caused by approximate optimization. Identical
        coefficients and monotonic decrease of the raw loss are not required.
\end{itemize}

\subsection{Task 5 (Optional)}

\begin{itemize}
  \item Reload T2 and make a new stratified 80\%/20\% split with random state 42
        (160 training and 40 test samples). Do not reuse the model or selected
        parameters from Task 2, whose search used all T2 samples.

  \item Use a StandardScaler--RBF SVC Pipeline with $C=1$ and
        $\gamma\in\{0.01,0.1,1,10,100\}$. On the training subset only, use
        shuffled, stratified five-fold CV with random state 42. Report mean fold
        training accuracy, mean validation accuracy and validation standard
        deviation. Plot the training and validation curves against logarithmic
        $\gamma$, with validation variability. Also fit each candidate on the
        full training subset to report its number of support vectors.

  \item Select $\gamma$ solely by the highest mean validation accuracy; ties use
        the first candidate in the listed order. Freeze the choice and evaluate
        only the selected model once on the independent test subset.

  \item Show the smallest, selected and largest $\gamma$ boundaries using
        training data. Analyze underfitting, overfitting, the training/validation
        gap and support-vector counts without assuming monotonic results.
        Discuss the uncertainty of a 40-sample test set: one changed prediction
        changes its accuracy by 2.5 percentage points. Leave
        \texttt{RUN\_TASK5=False} if you skip this task.
\end{itemize}

\section{Requirements}

\begin{itemize}
  \item Implement learning algorithms with \textbf{Python}. Basic data analysis
        libraries such as \textbf{NumPy}, \textbf{Pandas}, \textbf{Matplotlib},
        etc., are allowed to use. Tasks 1--3 and 5 may use
        \textbf{\textit{sklearn}}. In Task 4, the training algorithm, loss,
        gradients and predictions must use \textbf{NumPy only}; sklearn is
        allowed solely for the comparison model.

  \item Record the procedure of experiments in \textbf{Jupyter Notebook}.
        Use \textbf{Markdown} to explain your methods and results in English.
        Include the requested figures and tables with labelled axes and captions.

  \item Restart the kernel and run the completed notebook sequentially before
        submission. Record package versions, random seeds and CPU execution
        times. The provided settings use two search workers and 160-by-160
        visualization meshes; no GPU is required.

  \item Your submission includes:
        \begin{itemize}
          \item source code: the completed Python notebook (\textit{.ipynb})
                with executed outputs;
          \item \textit{report.pdf} file: your English report. Refer to
                \textit{Template.doc} for the format, and clearly state whether
                you completed Task 5.
        \end{itemize}

  \item The deadline is \textbf{November 11, 2026}.
\end{itemize}

\end{document}
